% 分野別類題: 最小二乗法（多項式フィット） 解答例
% コンパイル: TEXINPUTS="../../../00_common:$TEXINPUTS" latexmk -xelatex answers.tex
\documentclass[a4paper,11pt]{article}
\usepackage{preamble}
\usepackage{shoakubox}

\begin{document}

\kamokumei{類題演習 解答例 --- 最小二乗法（多項式フィット）}

\daimon{【解法】最小二乗法の解き方と導出}

\noindent
データ点 $(x_i, y_i)$ ($i=1,\dots,n$) に対して、モデル関数 $f(x)$（未知パラメータを含む）を当てはめる最小二乗法では、誤差二乗和
\[
S = \sum_{i=1}^{n} \bigl(y_i - f(x_i)\bigr)^2
\]
を最小にするようにパラメータを定める。$f$ が未知パラメータについて線形（例えば $f(x) = a x + b$ や $f(x) = ax^2+bx+c$）である場合、$S$ は各パラメータの2次関数であるから、最小値では各パラメータに関する偏微分がすべて $0$ になる。

\textbf{（1）直線フィット $f(x) = ax + b$ の場合。}
\[
S(a,b) = \sum_{i=1}^n (y_i - ax_i - b)^2
\]
とおくと、
\[
\frac{\partial S}{\partial a} = -2\sum_i x_i(y_i - ax_i - b) = 0, \qquad
\frac{\partial S}{\partial b} = -2\sum_i (y_i - ax_i - b) = 0
\]
これを整理すると、正規方程式
\[
\left\{
\begin{aligned}
a\sum_i x_i^2 + b \sum_i x_i &= \sum_i x_i y_i \\
a\sum_i x_i + b\, n &= \sum_i y_i
\end{aligned}
\right.
\]
を得る。これは $a, b$ に関する連立1次方程式であり、係数行列の行列式が $0$ でない限り（データが1点に退化していない限り）一意に解ける。

\textbf{（2）2次関数フィット $f(x) = ax^2+bx+c$ の場合。} 同様に $S(a,b,c) = \sum_i (y_i - ax_i^2 - bx_i - c)^2$ を $a,b,c$ で偏微分して $0$ とおくと、
\[
\left\{
\begin{aligned}
a\sum_i x_i^4 + b\sum_i x_i^3 + c\sum_i x_i^2 &= \sum_i x_i^2 y_i \\
a\sum_i x_i^3 + b\sum_i x_i^2 + c\sum_i x_i &= \sum_i x_i y_i \\
a\sum_i x_i^2 + b\sum_i x_i + c\, n &= \sum_i y_i
\end{aligned}
\right.
\]
という3元連立1次方程式（正規方程式）が得られる。

\textbf{（3）行列表記。} 一般に $f(x) = \sum_{j=0}^{m} c_j x^j$（$m$ 次多項式）のフィットは、計画行列 $A \in \mathbb{R}^{n\times(m+1)}$、パラメータベクトル $\bm{c}$、観測ベクトル $\bm{y}$ を用いて
\[
\|\bm{y} - A\bm{c}\|^2 \ \longrightarrow\ \text{最小化}
\]
という問題に帰着する。$S(\bm{c}) = (\bm{y}-A\bm{c})^{\top}(\bm{y}-A\bm{c})$ を $\bm{c}$ で微分して $0$ とおくと
\[
\nabla_{\bm{c}} S = -2A^{\top}(\bm{y} - A\bm{c}) = \bm{0}
\quad\Longrightarrow\quad
A^{\top}A\,\bm{c} = A^{\top}\bm{y}
\]
という正規方程式を得る。$A^{\top}A$ が正則であれば $\bm{c} = (A^{\top}A)^{-1}A^{\top}\bm{y}$ と一意に解ける。直線フィット・2次関数フィットの正規方程式は、いずれもこの行列形式の成分ごとの表示に他ならない。

\clearpage
\begin{mondaiboxS}{問1}
4個のデータ点 $(0,1),(1,2),(2,2),(3,4)$ に対し、直線 $y=ax+b$ を最小二乗法でフィットせよ。
\end{mondaiboxS}
\kaitou
$n=4$ とし、必要な和を計算すると
\[
\sum x_i = 6,\quad \sum x_i^2 = 14,\quad \sum y_i = 9,\quad \sum x_i y_i = 18
\]
正規方程式は
\[
14a + 6b = 18, \qquad 6a + 4b = 9
\]
第1式を $2$ で割って $7a+3b=9$。これを $4$倍、第2式を $3$倍すると
\[
28a+12b=36, \qquad 18a+12b=27
\]
辺々引いて $10a = 9$、すなわち $a = \dfrac{9}{10}$。これを $6a+4b=9$ に代入すると
\[
6\cdot\frac{9}{10} + 4b = 9 \ \Longrightarrow\ 4b = 9 - \frac{27}{5} = \frac{18}{5} \ \Longrightarrow\ b = \frac{9}{10}
\]
したがって求める直線は
\[
\boxed{\; y = \frac{9}{10}x + \frac{9}{10} \;}
\]
（検算: 残差 $r_i = y_i - (0.9x_i+0.9)$ は $r_1=0.1,\ r_2=0.2,\ r_3=-0.7,\ r_4=0.4$ で、$\sum r_i = 0$ かつ $\sum x_i r_i = 0\cdot0.1+1\cdot0.2+2\cdot(-0.7)+3\cdot0.4=0.2-1.4+1.2=0$ となり、正規方程式を満たすことが確認できる。）

\clearpage
\begin{mondaiboxS}{問2}
5個のデータ点 $(-2,4),(-1,1),(0,0),(1,1),(2,6)$ に対し、2次関数 $y=ax^2+bx+c$ を最小二乗法でフィットせよ。
\end{mondaiboxS}
\kaitou
$n=5$ とし、必要な和を計算する。$x_i$ が原点対称なので $\sum x_i = 0$, $\sum x_i^3 = 0$ となることに注意すると
\[
\sum x_i^2 = 10,\quad \sum x_i^4 = 34,\quad \sum y_i = 12,\quad \sum x_i y_i = 4,\quad \sum x_i^2 y_i = 42
\]
正規方程式は
\[
34a + 10c = 42, \qquad 10b = 4, \qquad 10a + 5c = 12
\]
第2式より直ちに $b = \dfrac{2}{5}$。第3式より $c = \dfrac{12 - 10a}{5}$ を第1式に代入すると
\[
34a + 10\cdot\frac{12-10a}{5} = 42 \ \Longrightarrow\ 34a + 24 - 20a = 42 \ \Longrightarrow\ 14a = 18 \ \Longrightarrow\ a = \frac{9}{7}
\]
このとき
\[
c = \frac{12 - 10\cdot\frac{9}{7}}{5} = \frac{\frac{84-90}{7}}{5} = \frac{-6/7}{5} = -\frac{6}{35}
\]
したがって求める2次関数は
\[
\boxed{\; y = \frac{9}{7}x^2 + \frac{2}{5}x - \frac{6}{35} \;}
\]
（検算: 予測値は $\hat y(-2)=\frac{146}{35},\ \hat y(-1)=\frac{25}{35},\ \hat y(0)=-\frac{6}{35},\ \hat y(1)=\frac{53}{35},\ \hat y(2)=\frac{202}{35}$ となり、残差 $r_i=y_i-\hat y(x_i)$ は $\sum r_i = 0$, $\sum x_i r_i = 0$, $\sum x_i^2 r_i = 0$ を満たす（正規方程式の各式に対応）。）

\clearpage
\begin{mondaiboxS}{問3}
データ点 $(x_i,y_i)$ ($i=1,\dots,n$) に直線 $y=ax+b$ をフィットする最小二乗問題を、計画行列 $A$ を用いて $\|\bm{y}-A\bm{x}\|^2$ の最小化として定式化し、正規方程式 $A^{\top}A\,\bm{x}=A^{\top}\bm{y}$ を導け。さらに問1の正規方程式と一致することを確認せよ。
\end{mondaiboxS}
\kaitou
計画行列とパラメータ・観測ベクトルを
\[
A = \begin{pmatrix} x_1 & 1 \\ x_2 & 1 \\ \vdots & \vdots \\ x_n & 1 \end{pmatrix}, \qquad
\bm{x} = \begin{pmatrix} a \\ b \end{pmatrix}, \qquad
\bm{y} = \begin{pmatrix} y_1 \\ \vdots \\ y_n \end{pmatrix}
\]
とおくと、$i$ 番目の行成分により $(A\bm{x})_i = ax_i + b$ となるから、誤差二乗和は
\[
S(\bm{x}) = \|\bm{y} - A\bm{x}\|^2 = (\bm{y}-A\bm{x})^{\top}(\bm{y}-A\bm{x})
= \bm{y}^{\top}\bm{y} - 2\bm{x}^{\top}A^{\top}\bm{y} + \bm{x}^{\top}A^{\top}A\bm{x}
\]
と展開できる。これを $\bm{x}$ で微分すると
\[
\nabla_{\bm{x}} S = -2A^{\top}\bm{y} + 2A^{\top}A\,\bm{x}
\]
であり、これを $\bm{0}$ とおくことで正規方程式
\[
\boxed{\; A^{\top}A\,\bm{x} = A^{\top}\bm{y} \;}
\]
を得る。ここで
\[
A^{\top}A = \begin{pmatrix} \sum x_i^2 & \sum x_i \\ \sum x_i & n \end{pmatrix}, \qquad
A^{\top}\bm{y} = \begin{pmatrix} \sum x_i y_i \\ \sum y_i \end{pmatrix}
\]
であるから、成分ごとに書き下すと
\[
\sum x_i^2\, a + \sum x_i\, b = \sum x_i y_i, \qquad \sum x_i\, a + n\, b = \sum y_i
\]
となり、これは問1で偏微分により導出した正規方程式
\[
14a+6b=18, \qquad 6a+4b=9
\]
（$\sum x_i^2=14,\ \sum x_i=6,\ n=4,\ \sum x_iy_i=18,\ \sum y_i=9$ を代入したもの）と完全に一致する。（検算: 問1で得た解 $a=b=\dfrac{9}{10}$ を $A^{\top}A\bm{x}$ に代入すると $A^{\top}A\bm{x} = \begin{pmatrix}14\cdot\frac{9}{10}+6\cdot\frac{9}{10}\\ 6\cdot\frac{9}{10}+4\cdot\frac{9}{10}\end{pmatrix} = \begin{pmatrix}18\\9\end{pmatrix} = A^{\top}\bm{y}$ となり、確かに正規方程式を満たす。）

\end{document}
